\section{Magnetic field control}

The magnetic field control system consists of the magnetic shield, the set of
shim coils, and the internal and external magnetometers. Together, these
components allow for interaction region operation with low levels of ambient
magnetic background present. This section will discuss the design requirements
of the system, present the physical details of its implementation, as well as
the performance achieved with the present shielding system.

While the design, building, and characterization of the shield can be done in a
straightforward manner by emulating past work, the underlying science is
frequently not obvious and is often left out of the literature on the topic.
This is understandable, since understanding how the shield works is rarely of
interest to the experimentalist desirous of getting an atomic or molecular beam
operational. In what follows, I will also have to leave several lines of
exploration unfinished as they are not immediately required in order to proceed
with the \CENTREX\ experiment.  However, it may be a starting point for future
work requiring higher performance.

\subsection{Design requirements}


\subsection{Theory and related work}


\subsection{Physical structure}

The magnetic shield con s i s t s of four conce ntr i c cyl i ndr i ca l lay.ers
with open ends. Three of these I built by wrapping metglase foil around PVC
pipes, while the fourth, i nnermos t layer has been welded out of two sheets of
mu-metal rolled into a t ubul ar. form, and aneeale'- Themu-metal tube, bei ng 1
mrn thick, is r elat i vel y fra� gile; in Ord tr< to l i mt.t its motion, it is
sa ndwi c he d bet wee , the bard pl ast i c PVC tubes of the i ner metglas
layer and that of the shim coils. Thei nner twometglas layers are spaced apart
using small-diameter PVC tubing; the :same m thod is employed to suspend the
glass vacuum chamber of the i nt er act i on region cte n t er e d inside the
shim coils tube. The outermost tube :is mounted on a n 80/20 fra.ijle usi ng a
set of seven al umi num s emi - c ri c u l ar plates

\paragraph{Layers, spacers, mounting}

The three metglas layers are all built in the same way. Piret, the plasticr tube
( provided by the Harrison Machine \&\ Plastic Corp.) is cut to size with a
handsaw if n ecessary ( see Table for layer lengths and material details) and
roughly clea1ed w. i t h water and paper towels on the outside a nd as far as
possible on the inside to .remove the dirt it comes with, probably due to
outdoors storage. TheR, I cut 8x12 pieces ofll metglas (7x1 for the 16 11 layer
==:::::FACT CHECK) to the length reouiredto wrap around the tube.with approx. 1
11 of overlap. The pieces are secured to the PVC tube and to each other with a
few short pieces of Kapton tape (McMastar p.n. xx..<L{) so that they do ri ot
fall down during assembly.

The prepared and cleaned tubes are put on a stand with wheels unoer them so they
can be rotated, fa ilitating the circumferetial wrapping of the metglas fo,111:.
The first piece of tape is wrapped ab9ut an inch front:the e11d . of t,he tube;
leaving this much be.re PVC . help with handling the approxJ 100-200 lb tubes.
The second piece is wrapped so as to have about an inch of overlap with the
first, ancl the remaining six pieces are attached a nalogously. Once the first
layer is done, the .wrapping resumes at the position of the first metgl ss
piece. This procedure is repeated for a total of 12 layers total. (Due to axial
and circumfere tial overlaps, the true thickness of the shield is not perfectly
homogerEous, but in practice we find that this is unimportant aslong as the
overall shieldi ,g is sufficien high.)

To secure the metglas pieces and prevet them fran c'.oming loose in case the
kapton ,ta pe fails, the entire tube is wrapped in plastic foil {McM p.n.
AXXXA). Following that, there are .12 turns of degaussing wire wrapped in a
toroidal fashioi around the tube. The wire 1s 22 o;r 24 gaug lacquered copper
wire (McM p.n .JUX), sufficientyl thick for approx.  2 amps of degaussing
current. The 12 turns are uniformlv distributed around the pirc mfere ce of the
tube; we qelieve thism may be - important for the performance of the degaussing.
On the outennost ( 20") tube only, we also have a solenoidal coit with approx 1
turn/in spacing for rough caneellation of axial fields.- Given the metglas
wrapping geometry, the solenoidal coil would not be effective for aegaussing.)
Finally, :tM::::eai each tube is again wrapped in plastic wrap to secure the e:e
degaussing and ( for L011 tube Otl y) the solenoidalcoil


The mu=metal layer doe8 not,' of course, require wre.ppinc?: with metglas,
oil.ly the degauss ing coil aud outer plas ic foil. However, Cere must be taken
to reinforce the insulation of the degaussing wire at the ends of the cylinder;
the sharp edges oan easily strip the magnet wire coating a d snort the tur.,s of
the wire to the conductive cylinder, which does not happen with thePVC tubes. I
used shrit1k wrap to do this, but it required cutting the wire a.t each turn a.t
'ie both ends, adding theshrink tubes, soldering the wires back together and
heatigg tfie ihsulation to 11, make it adhere. It would have been be ter to use
a wire with more durable insulation instead.

\paragraph{Coils, current source}

The set of shim coils co of 16 segments three diaensions each(.;61 The two
transverse dim (Bx and By) are plemented as pairs of one-loop saddle coil ,. 1le
the axial fiel if controlled by simple circular current loop. This gives a tota
of 5x16a80 degrees of freedan. The 22g magnet wires are fixe n place by being
glued into a series of grooves on the surface o t 1411 PVC tube using neutral-
cure RTV silicone sealant (Do 37). e sise of the tube make111 it impractical to
use sta da.i:d milling machines; the circular grooves can easily be cut using a
har saw, while the axial ones were done on a custom-ma.de setup involving a hand
drill clamped on a 0/20 frame, which in tur is mounted using linear bearings to
a larger frame. The setup is not as rigid as a true milling machina would be,
which is why it is importa t to use a constant force whilemoving the carriage
forward to avoid cutting wavy grooves. A used a system of ropes, pulley , and
weights to achieve satisfactory r s u.t ,s .

With 80 different pair of wires (twisted into a twisted pair using a hand
drill), cable management is essetial so as not to confuse the coils with each
other. For each coil,a linear groove is cut fr-om one end of thetube to
thelocation of the coil; thus all coils can be connected to at one end of the
shield. There, the sets of wire pairs belonging to Oile dimension are bundle'
together into five bundles of 16 pairs using zip ties; each wire pair is
temporarily labelled by a numbered piece of kapton tape at the end. For ease of.
connecting a d <iisconecting, these wires are then attachea to adapterPCBs that
are permanently attached to the 80/20 support frame of the shield.The adaptor
boards bundle four pairs of wires into on RJ45 jack to enable using ether et
patch cables for ease of connection.


\paragraph{Support structure}

The considerable total weight of the magnetic shield (900 lbs of plastic, plus
of metglas, XXlbs of mu-metal, and other materials) requires a sufficiently
atroM support structure. I our case, it is assembledfrom 311 'd0/20 al uminum
rails ( profile .umber .{A,{) and placed on wheels {McM p.r1. AXX) so the setup
ca.n in principle be moved .by one person from the lab upstairs, where ,ii -a
the setup was built, to theone downstairs where it will be used. The
supportstructure consists of three 4' segments mounted togethwer to form a
rectagular box. Bails going across the top of 'the box hold alurnnium plates (
1/,4" thick) with semi-circular cutouts. To keep the plates from damaging the
delicate outer mei;glas layer/ the shar corners of the cutouts are protected
with hard ylon tppe (AAu wide, U 11 thick, McM ,p. 1 . ,1 ).


\subsection{Operation and measurements}

The cha.ra cterizationand use of the shield is an iterative process requiring a
rppeated seauenc.e of three opera't ions: magnetic.field mea= surement,
degaussihg, and shimming.

I distingu sh twotypes of mea.surment of the gue t i c field of interest to us.
Residual field is hhefield inside the shield assuming changes in the extern l
magnetic condi t i uns . In our four-layer shieid, this is the domi na nt mag ne
t i c field inside the interaction region arxi is driven primarily by the i nt
eract i on between the me.g etization of the f'er oma g i1e t i <:,; shield
layers with the exteral mag ne t i c field. Despite our best at t emp t s at
degaussing, theser esi du b.l f i el d s i are typically: of hhe. oroer of 50-
200 uG- in the  er of the shield 

\paragraph{Field measurement}

During shield characterization, the primary measurement is of the field in th
central region of the shield. During the 4se of the shield in a
beam-line.experiment, this region is difficult to access; it is also the region
of greatest interest since all molecules included in the exp= er1me t pass
through it. D uring characterization 1-m ea attached a 411x411 alumnium
rectangular tube { i side dime siuns A.i(xll") by suspending it from them shield
mounting structure at each ene . I used a system of pulleys to pull the se sor
carriage using a stepper -motor( and a rot� ary e coder to measure the position
of the se sor al ng the length of tube). The sesor itself ( Bartington
Ma.g03MS-100) is a 3-cha,m el flux= gate magnetometer with a measuring range of
1G a nd a typical elEctroic offset of 9'l several hundred microgauss.( The
spQ,cif'ied noise level of �.XX. uG/rtHz, integrated over the a-w at-a-ef --the
s e 3 kHz bandwidth of the sensor, is below the mechanical stability of the
sensor trai1sla� tion setup with a typical repeatability of approx 2uG. FACT
CHECK THIS).

The electronic offset of the sensor and its associated electronics is
significant clt is specified as XXX uGi, with a temperaturecoefficient of XXX..
uG/c; whil in practice values of xXX. uG are typical. The offset calibration is
done by taking a measurement and immediately repeati ng it with the sesor
rotated by 90 deg. This- is accomplished by mounting the sensor in a short
rectangular aluminum tube with enough plastic tape on the outside so that these
sor carriage fits tightlyinside the long soua.re tube going through the center
of the shield. After a n.easu� rement is done, the sesor ls pulled out of tl:le
l ong tube and reinserted rotated upside down. The accuracy of the 4 v offset
calibration then depe .ds on how well-c1ettered the sesor is in the carriage,
how much the field inside the shield changes during the time it takes to perform
the c. libration, how much theelectroic offset itself changes in the same time,
ard how pe . a s quare the cross-sectio s of the long tube and the carriagesare.
ather than attempt to estimate the ma.gnitm es of these effects i idually, I
estimate the precision/accruacy of the calibration fromthe calibration curve
itself; see Fig. AU for a typical exa.rnple. It shows the two travsverse
compuI1ents of the magnetic field ( Bx and By) a-s measured by the sensor in its
nonnal orientation,as well as with the sensor inverted. The data is presented ae
in the symmetrized form: tiormal+inverted)/21- which corresponds to the se11s or
offset, and (normal-inverted)/2, which corresponds to the physical magnetic
field plus any calibration imperfections. Given the tim scaleof the changes in
the electronic offset is primarily driven by temperature changes, which are
negligiblein the time it takes to do the calibration (around 2 min), we expect
the curves showing the offset to be perectly flat within the noise of the
measurement, but that is typically not the case. Therms deviations from the
averagevalue I identify as the error in the offset calibration, cosisting of all
the aforementi m ed fact rs.

While one may take the(normal-inverted)/2 superposition s the best measurement
of the physicL fields, doing so means that tne error is different at each
location along the shield. Since we expect the offset to be essentially constant
during the calibratino, I believe it w a better practice to estimate theoffset
from the averagevalue .of the+ superposition, and then simply -subtra.ct this co
stant value from all subsequelt measu.ramets. This way, the overall offset is
subject to a +/- XX uG uncertainty (typ.), bu the shape of the field pattern is
known much more precisely.The typical repeatability of a measurement is within
M%, as i M a s nown in an ex8 rnple measurement i r1 Fig. UK.

In addition to the ra dom errors i .duced oy the iJ,echa. 1.cal �""'., ,
,..,..d.,(>.., J inatabil ty of th.e measuremen.t setup we have to pay Jte
i"tr'onto the' sys tematic errors induced by thesignal chain from the sensor
onwards. The sensor a,1d the preamp have a limited bandwidth; 3 kHz for the se
sor, and a settable low-pass filter for the preamp ( between 0.1, 1, 10, 100,
XL{ Hz). Thia must not be set too low in relation to the speed of the sensor
motion. In practice, I use the lowest...reasonable LPP cutqff freq., verifying
my choice by repeating the mea s ur eme flt with the nex t higher avai l ab1'e
freaue,cy:.. If the two me a s ur eme nt are s uf f i c i e nt l y close to each
oth er , I use the lower of the two f'reqs. use a similar process to decide upon
the d1gital LPF inherent to theADC ( ADxx,xx.. used in a hoxoo� bui ot ADC
board)

\paragraph{Dynamic field measurement}

On the other hand, the atte uat i on of t he ex ter ,a l :(iel d s i s ,+-1;,Q
11ot p e rf ec t . If an infinitely l a.rge He lm h ol z pa ir we r e to drive
a perf ect l y unif orm..and s i usoidally varyi ng ma g n e t i c field in the
region of space occupied by the magnetic shi el d , t h en a smaller sinusoidal
var iat i on of t hef i el d inside the sni eld could be det ect ed that is of
the same freque cy as the e.;Ct e r 11a l field. In practice, there are
practical limits to the size of exter al coils. In one vre s i on of' t hemea s
ur eme n t , I have attached a room-sized loop of wire to t he walls of the lab
and dri ve11 them witha function generator while eo N pavsing the ma g ne t omet
er re ding through a lock-in amplifier. Re pe a t i ng t h e mea s ur e me n t
irn the a absence of any shielding gi ves the mag itude of t ne external field�
variation. Di v i d i ng t h i s by the mag itude measured with the sni el d in
place gives the shieicting factor. SeeFig. XXX for seve al variations of this
measurement for va ni ous s hi el d i g c onf i gur at i ons.. C

Thedsadvatage of this method is that the magnetic shi el d has to be� removed f
to dete rmi ne the ex t e r na l ma g net i c field, since the shape of the
exteral coils is difficult to easure preci sel y. To find theSF o the final
shield, I have used smaller rectagular coils of well-known dime nsi ons cent e r
e d around the mag ne t i c s hi el d ; the field they apply, while not uni f
orm,, cau be calculated using a closect-f orm expression ( see eqn )(XX. in pa
per YYY ). See Fig. XXX for details about EJ,e recta gular coilo-Z:: om <r t ry
. The results of the meqsurements in this configuration are discussed
insect:performance.

\paragraph{Degaussing}

Degaussing is a method of reducing the residual field inside a magnetic shield
baaed on the following observet.ion:the field inside a hollow cyl1nder made of a
ferromagnetic IDB.terial,when placed in a region of a given external field, can
be greatly reduced by driving asinusoidal magnetic field with a gradually
decreasing e ,velope.see.Fig. XA for an-example of the degaussj..ng.. g ometry,w
a v.e :lorm, ard the result improvement in the shielding.

The hardware used...for degaussing in this experiment is as follow.s Each of bhe
four shlelding ayers H has its own 12 - turn degaussingcoil. These are connected
to a PCB with a set of relays that can connect each of them to the degaussing
current, but only one at a t e. During degaussing,the outer ' ayer isalways
cormectetl- first,followed by the secorx:l outennost layer ollowed by the
innermost metglas layer, iollowed by the mu-metallayer. i3ce the residual field
is primarily determined by the innemost layer, that one is degaussedlast lest it
be magnetized by accident duri g the degaussing of the other layers. The relays
PCB is in turn co nnected to an isolation transformer for the purpos,e of DC
blocking We have considered alter nate forms of DC blocking, iuch as a capacitor
e11-R-J., but the requ ired capacitance a� -aQ88":iei --a ae 1 was too high at
the low frequetlcies typica.llyusedin degaussing ( down to 1 Hz). The
transformer is driven by a power amplifier,Kepco BOP36-12M), which in turn is
driven by a home-built DjC boa.rd. S ee Fig. XXX.. for a schema.tic of the
degaussing wirig connections.

A wide rage of degaussing waveforms have been used in the literature for
exampee, li ar or polyllomial enveloves {ref. XU) or envelopes ' specifically
optimized for a linear decrease in themagnetization{?) of the ferromagnetic
material. In this experiment, I use a finite umber of a.a s inusoidaldegaussil18
pulses. The puls s can Optionally be separated by and adjustable delay to probe
tfie evolution of the residual field from pulse to pulse. Each subsequent pulse
has its amplitude decreased by a constant factor comared.with the preceding
pulse; thus our protocolw ia similar to, but identical, using an exponentially
decreasing e velope multiplyig a sine wave as used in refai XXX and YYY, as well
as the ACME ex eriment {ref. ZZZ) .Pig. XXX identifies the parameters of the
dega� ssing waveform graphically

Despite spedi1g aig,,ificant amounts of time o optimizi g the degaussing
parameters, I am only able toe claim the following. tega.ussing freque(1Cy ha.s
to be slow enough that the skin effect does not significantlyaffe ct the
degaussi11g1; the initial amplitude has to be large enough to drive the shield i
1to saturation; the -final amp. has to-b e-small enough that tne shieI -has
recovered from saturation and is circlin-g near a point around the origi of the
hysteresiscurve; the number of degaussing pulses has to be large e ou enough so
that the decrease in amplitu is gradual.;'.l'.he exact values of the parameters
metioned f n tae prece'1ing ta'tements need to be found experimnet lly for the
shield geometry used; seJ-Fig. for the data correspondi to my
optimizationattempts. The final jtalues of he degaussnig parameters used are
also specified in:hat/ figure; these are not necessarily optimal, but the
procedure works well enough for our purpose.

Exp1-.ora tion of these paramters is further comp ic ?.ted by he fact that there
are four layers to deg,uss, ard degaussing one moy inadv erted ty m gnetize the
others-. I have tried degaussing sevral or all layers in p:rallel, but have not
found any s ignific nt improv mer:it and ti:ius h ve not further explored that
method., However, while building the indivictual metglas layers, I have
measured the magne izationof each layer before a nd after degaussing; see Fig.
X.U for this data.

Given the 1'lportance of the long-tenn sta.bility of the mag ot . le field
inside the shield, it is interest{ilg to explore how stable th degaussing
process is, i.e., how much does the field change when degaussing., repeatedly,
or MW- what happens when the field has driftedaway from the oeg,ussedvalue and
the shieldis degeusae.d again. (Make a figure with the answer to this question.)
The inter ctiofl of degaussing with the shimming coils is explored in
the.followgi section.1

\paragraph{Shimming coils}

Once the shield has been degaussde as well as is possible, the rema.iic.g
magnetic field has to be Ca.llcelled using the systern of shim coils. The
accuracy of this depends on the measuremer t of the residual field,- as well as
on the coils themselves. Duri1 shield characterization, we have access to the
center of the shield, and can mea.s ure hte fields t within 1.:-2 uG. These can
be shimmed away to withinX uGrms; the limit is set by the1umber of coil segments
and their size.

The shimmmng process consists of t o steps. First, the individual fitld patterns
generated by each of the 80 shim coils have to be measurceti. I do this by
Jllowing a constant current thr@ugh thecoil under test, measuring t the magnetic
field, then reversing the direction of the current ard re� pec�ing the
measurement. By subtracting the .two measuremnets.,applyi1g smoothing (
usingsavglofilterl(l., XX , ll)), andlf masking the parts of the curve where the
coilfiela is expected o-4-- be egligible(see Fig.  XXX.for an exampee of
thecalibration curve). Once the calibration curves for all the coils are known,
I, usea simple least-squares mi:1imization routine to determine the optimal
currents to use .so as to mi i-. ze the magnitude of tne residual field insi'de
the shield. While I could limit the opti ization routine to oly minimize the
field along the f rtical axis, this being thecomponent we are most sensitiveto
in themolecular =beam experiment, I do, not anticipatea significant improvemnet
in the fila-ielein shimmmng by doing so give.1 th8fl three sets of coils ( two
pairs of transverse coils, ane set of axial) produce fieldsthat are mqstly- -
rthogonal to,each other. The shimmed residual field is shown in Fig. XXX ion
sect:performance.

\subsection{Performance and limitations}

The design requireme ts for the shieldspecified a 10 uG maximum averag e re
du.a,l field. This ha een met/ t; see Fig. XXX for a ttypical measurement of the
fields inside after shimming.

-Of pa.rticuler i nterest is how the fields de ve l op after the shimming
procedure. To this end, I have followed their evolution ovre the course of three
weeks, finding that nei t he r the average nor the ma Va l ue of any of the two
transverse compil e nt s d owe not exceed the specified 10 uG ( see Pi.g. Xll).
Thedata in that plot has been obta,i e d by mea� suring the fields evecy hour
while keeping all the controllable para� meters of the shield constat. Unf or t
una t el y, at the time of the mea� sureme t , I did not have avaiable the me a
ri s to monitor the magnetic field ext erna l to the magnetic shield, or at a y
other point insi.e the shield except on-axis.

Sinc e <iegausaing is used to sig if i cant l y reduce the i nt er r1a l
magnetic f i eld,, it is of interest to observe whether dega us s i ng could be
used per i odi c til l y to restore the initial low ield coi1'di tions even
after shimming the fields. Unfortunately, 'it appears that degaussing leaves the
shield in a state of a small but signi f i ca t1t internal .:field-. Since the
shield is stable eough, I believe better performa ce is obt tai ned by
degaussing only at the beginning of mea s ur eme nt s , then shimming the
shield, ad letting it stay hat way.

It is true that just because the shield stayed in a low-field condition for
three weeks once, it is not certain it will al wa-y s stay that way for as l
ong-. In o:t;her words, it would beuseful t.o comeup with a model for predicting
how l ong- t he shield ,st aya this way, or what the pr oba bi l i t y - i s
that is has exceeded lO -uG. If we model ,t he average i nt er nal magne t i c
field as a radom walk, t hen -we can associate a diff usi on c ons t a nt to
this process and calculate- the probabilities. {Include the calculations for
radom walk.) Of course, there is not physical reason why the growth in interal
magnetic field should be a ra ndomw al ; it would be of interest to compare
other possible stochastic models for this behaviour.

In addi t i on to the aQ,ee e static inter ,a l field, it i,s also important to
ensure the exter al mag net i c -f i el d noi se is suf f i c i ent l y a t t
e-iua t ed inside. My measurements of the dynamic shielding fac-0tr placeit
around 30, 00-0 at DC ( see F i g . XXL), which sh oul d . reduc-e the expected
1mG varia� tions in lab field well below the l evel at which they would impact
the molecular-beam mea s ur eme nt s . The frequency depend ec e of the dyna mc
SF shown i nt ha t figure shwo s an i ncr ea se i n shielding as f requency
incre� ases, which ma ke 1;; s e nse in light cf skin ef1'ect cons i der a t i
ons . For mu� metal of 1 mm thickess and mur = 50k, a frequency of XX Hz would
cor� respond to a ski n depth comparable to the i nner - l aye r t hi c k11 e e
s . (The decrease in shi eldi11g once tne maximum is achieved is l i kel y not
phys i c pl , a na reflects som limitation of my mea sur i n setup. I have, howe
ver ,  been so far unable to find out shat causes this.)

The spatial epena e nce of the dynamic SF is as shown in Fig. Xli. It w.-ill be
noted that the pattern is asymmetric, as if one side of thee� shield works
better than theother. That i s al so unl i ke l y to be the c@se.
ImageImageFuture work may establish whet h re from my todo 1 i st ) 


While the dynamic shieldi follows the approximate theoretical model, the static
shielding does not.b It is not clear why aegausaing achieves such good shieldi ,
nor is it known wti,t limits it from working even better. It appears clear from
a simple model that degaussing magnetizes the shield in such a way as to oppose
the external magnetic fields, that being the configuration of lowest total e
argy
